\chapter{Hall Subgroups}

\providecommand{\Sym}{\operatorname{Sym}}
\providecommand{\PSL}{\operatorname{PSL}}
\providecommand{\Syl}{\operatorname{Syl}}
\providecommand{\Hall}{\operatorname{Hall}}
\providecommand{\Dih}{\operatorname{Dih}}
\providecommand{\Op}[2]{\mathrm{O}_{#1}(#2)}
\providecommand{\CG}[2]{\mathrm{C}_{#1}(#2)}
\renewcommand{\NG}[2]{\mathrm{N}_{#1}(#2)}

\section{Definition and Basic Examples}

\begin{definition}[Hall subgroup]
Let \(G\) be a finite group, and write
\[
    |G|=p_1^{a_1}p_2^{a_2}\cdots p_t^{a_t},
\]
where \(p_1,\ldots,p_t\) are the distinct prime divisors of \(|G|\). Define
\[
    \pi(G)=\pi(|G|)=\{p_1,\ldots,p_t\}.
\]
For a subset \(\pi\subseteq \pi(G)\), a subgroup \(H\leq G\) is called a
\emph{Hall \(\pi\)-subgroup} of \(G\) if
\[
    \pi(H)=\pi
    \qquad\text{and}\qquad
    \gcd(|H|,[G:H])=1.
\]
Equivalently, \(H\) is a Hall \(\pi\)-subgroup if \(|H|\) is a \(\pi\)-number and
\([G:H]\) is a \(\pi'\)-number, where
\[
    \pi'=\pi(G)\setminus \pi.
\]
\end{definition}

\begin{remark}
If \(\pi=\{p\}\), then a Hall \(\pi\)-subgroup is precisely a Sylow
\(p\)-subgroup. Thus Hall subgroups generalize Sylow subgroups by allowing
several primes at once.
\end{remark}

\begin{topic}[Notation]
When a Hall \(\pi\)-subgroup has been chosen, we may denote it by
\[
    G_{\pi}.
\]
If \(\pi=\{p\}\), we often write \(G_p\) instead of \(G_{\{p\}}\). Also,
\[
    \pi'=\pi(G)\setminus \pi.
\]
This notation does not imply uniqueness: Hall subgroups need not exist, and even
when they exist, they need not be conjugate in arbitrary finite groups.
\end{topic}

\begin{example}[Hall subgroups in \(A_5\)]
Let
\[
    G=A_5,\qquad |G|=60=2^2\cdot 3\cdot 5.
\]
Then
\[
    \pi(G)=\{2,3,5\}.
\]

The point stabilizer of the natural action of \(A_5\) on five letters is isomorphic
to \(A_4\), and has order \(12=2^2\cdot 3\). Hence
\[
    A_4\leq A_5
\]
is a Hall \(\{2,3\}\)-subgroup.

A Sylow \(5\)-subgroup is cyclic of order \(5\), so
\[
    C_5
\]
is a Hall \(\{5\}\)-subgroup, equivalently a Hall \(\{2,3\}'\)-subgroup.

However, \(A_5\) has no Hall \(\{2,5\}\)-subgroup and no Hall
\(\{3,5\}\)-subgroup.

Indeed, suppose first that \(H\) is a Hall \(\{2,5\}\)-subgroup. Then
\[
    |H|=2^2\cdot 5=20,
    \qquad
    [A_5:H]=3.
\]
The left action of \(A_5\) on the set of left cosets \(A_5/H\) gives a
homomorphism
\[
    A_5\longrightarrow \Sym(A_5/H)\cong S_3.
\]
Since \(A_5\) is simple, the kernel is either \(1\) or \(A_5\). The action is
nontrivial, so the kernel is not \(A_5\). Hence the map would be injective,
which is impossible because
\[
    |A_5|=60>|S_3|=6.
\]
Therefore no Hall \(\{2,5\}\)-subgroup exists.

Similarly, if \(K\) were a Hall \(\{3,5\}\)-subgroup, then
\[
    |K|=3\cdot 5=15,
    \qquad
    [A_5:K]=4.
\]
The coset action would give a nontrivial homomorphism
\[
    A_5\longrightarrow S_4.
\]
Again simplicity forces this map to be injective, which is impossible since
\[
    |A_5|=60>|S_4|=24.
\]
Thus no Hall \(\{3,5\}\)-subgroup exists.
\end{example}

\begin{example}[Hall subgroups need not be conjugate]
Let
\[
    G=\PSL_2(11).
\]
Then \(G\) is simple and
\[
    |G|=660=2^2\cdot 3\cdot 5\cdot 11.
\]
This group has subgroups \(H\) and \(K\) such that
\[
    H\cong A_4,
    \qquad
    K\cong \Dih_{12},
\]
where \(D_{12}\) denotes the dihedral group of order \(12\). Both \(H\) and
\(K\) have order \(12=2^2\cdot 3\), so both are Hall \(\{2,3\}\)-subgroups of
\(G\).

However,
\[
    A_4\not\cong D_{12}.
\]
Therefore \(H\) and \(K\) cannot be conjugate in \(G\). This shows that Hall
subgroups are not necessarily conjugate in arbitrary finite groups.
\end{example}

\section{Hall's Theorem}

\begin{lemma}[The subgroup \(\Op{p}{X}\)]
Let \(X\) be a finite group and let \(p\) be a prime. Define
\[
    \Op{p}{X}
\]
to be the largest normal \(p\)-subgroup of \(X\). Equivalently,
\[
    \Op{p}{X}
\]
is the product of all normal \(p\)-subgroups of \(X\). It has the following
properties:
\begin{enumerate}[label=\textup{(\arabic*)}]
    \item \(\Op{p}{X}\trianglelefteq X\);
    \item \(\Op{p}{X}\) is a \(p\)-group;
    \item \(X/\Op{p}{X}\) has no nontrivial normal \(p\)-subgroups.
\end{enumerate}
\end{lemma}

\begin{proof}
The product of normal \(p\)-subgroups is again a normal \(p\)-subgroup, so the
product of all normal \(p\)-subgroups of \(X\) is the unique largest normal
\(p\)-subgroup of \(X\). This proves (1) and (2).

For (3), suppose that \(X/\Op{p}{X}\) has a nontrivial normal \(p\)-subgroup
\(N/\Op{p}{X}\). Then \(N\trianglelefteq X\), and \(N\) is a normal
\(p\)-subgroup of \(X\), strictly containing \(\Op{p}{X}\). This contradicts
the maximality of \(\Op{p}{X}\).
\end{proof}

\begin{lemma}[Minimal normal subgroups of solvable groups]
Let \(X\) be a finite solvable group. If \(1\neq N\trianglelefteq X\) is a
minimal normal subgroup of \(X\), then
\[
    N\cong (C_p)^d
\]
for some prime \(p\) and some integer \(d\geq 1\). In particular, \(N\) is an
elementary abelian \(p\)-group.
\end{lemma}

\begin{proof}
Since \(X\) is solvable, every subgroup and quotient involved in the derived
series is solvable. In particular, \(N\) is solvable. Since \(N\) is nontrivial
and normal in \(X\), the commutator subgroup \(N'\) is characteristic in \(N\)
and hence normal in \(X\). By minimality of \(N\), either
\[
    N'=1
    \qquad\text{or}\qquad
    N'=N.
\]
The second possibility is impossible because \(N\) is solvable. Therefore
\(N'=1\), so \(N\) is abelian.

Let \(p\) be a prime divisor of \(|N|\). The subgroup
\[
    N[p]=\{x\in N:x^p=1\}
\]
is nontrivial and characteristic in the abelian group \(N\), hence normal in
\(X\). By minimality, \(N[p]=N\). Therefore every nonidentity element of \(N\)
has order \(p\), and hence
\[
    N\cong (C_p)^d
\]
for some \(d\geq 1\).
\end{proof}

\begin{lemma}[Hall subgroups pass to normal subgroups and quotients]
Let \(N\trianglelefteq G\), and let \(H\) be a Hall \(\pi\)-subgroup of \(G\).
Then:
\begin{enumerate}[label=\textup{(\arabic*)}]
    \item \(H\cap N\) is a Hall \(\pi\)-subgroup of \(N\), after replacing
    \(\pi\) by \(\pi\cap \pi(N)\);
    \item \(HN/N\cong H/(H\cap N)\) is a Hall \(\pi\)-subgroup of \(G/N\),
    after replacing \(\pi\) by \(\pi\cap \pi(G/N)\).
\end{enumerate}
\end{lemma}

\begin{proof}
Since \(H\) is a \(\pi\)-group, \(H\cap N\) is also a \(\pi\)-group. Moreover,
\[
    [N:H\cap N]\mid [G:H],
\]
and \([G:H]\) is a \(\pi'\)-number. Hence \([N:H\cap N]\) is a \(\pi'\)-number,
so \(H\cap N\) is a Hall subgroup of \(N\) for the relevant set of primes.

Similarly,
\[
    HN/N\cong H/(H\cap N)
\]
is a \(\pi\)-group, and
\[
    [G/N:HN/N]=[G:HN]\mid [G:H],
\]
so this index is a \(\pi'\)-number. Therefore \(HN/N\) is a Hall subgroup of
\(G/N\).
\end{proof}

\begin{lemma}[A complement step]
Let \(G\) be a finite solvable group, and let
\[
    M\trianglelefteq G
\]
be a minimal normal subgroup. Suppose
\[
    M\cong (C_p)^d
\]
and
\[
    p\nmid |G/M|.
\]
Then \(M\) has a complement in \(G\). That is, there exists a subgroup
\(L\leq G\) such that
\[
    G=ML
    \qquad\text{and}\qquad
    M\cap L=1.
\]
Equivalently,
\[
    L\cong G/M.
\]
\end{lemma}

\begin{proof}
We argue by induction on \(|G|\). The statement is clear if \(G=M\).

Put
\[
    \overline{G}=G/M.
\]
Since \(\overline{G}\) is nontrivial and solvable, there exists a prime
\(q\in \pi(\overline{G})\) such that
\[
    \Op{q}{\overline{G}}\neq 1.
\]
Let
\[
    \overline{Q}=\Op{q}{\overline{G}},
\]
and let \(Y\) be the full preimage of \(\overline{Q}\) under the quotient map
\(G\to \overline{G}\). Then
\[
    M\trianglelefteq Y\trianglelefteq G,
    \qquad
    Y/M=\overline{Q}.
\]
Choose \(Q\in \Syl_q(Y)\). Since \(M\) is a \(p\)-group and \(p\neq q\), we have
\[
    Y=MQ
    \qquad\text{and}\qquad
    M\cap Q=1.
\]

We split into two cases.

\medskip

\noindent
\textbf{Case 1: \(\CG{M}{Q}\neq 1\).}

Since \(M\) is abelian,
\[
    \CG{M}{Q}=Z(Y)\cap M.
\]
Thus \(\CG{M}{Q}\) is the Sylow \(p\)-subgroup of \(Z(Y)\). Hence it is
characteristic in \(Z(Y)\), and \(Z(Y)\) is characteristic in \(Y\). Since
\(Y\trianglelefteq G\), it follows that
\[
    \CG{M}{Q}\trianglelefteq G.
\]
By minimality of \(M\),
\[
    \CG{M}{Q}=M.
\]
Hence \(Q\) centralizes all of \(M\), so
\[
    Y=M\times Q.
\]
In particular, \(Q\) is characteristic in \(Y\), and therefore
\[
    Q\trianglelefteq G.
\]
Now apply induction to \(G/Q\). The image \(MQ/Q\) is a minimal normal
\(p\)-subgroup of \(G/Q\), and
\[
    (G/Q)/(MQ/Q)\cong G/MQ
\]
has order not divisible by \(p\). Hence \(MQ/Q\) has a complement in \(G/Q\).
Taking the full preimage of this complement gives a complement to \(M\) in \(G\).

\medskip

\noindent
\textbf{Case 2: \(\CG{M}{Q}=1\).}

We first observe that
\[
    \NG{M}{Q}=\CG{M}{Q}=1.
\]
Indeed, if \(x\in M\) normalizes \(Q\), then for every \(q\in Q\),
\[
    xqx^{-1}q^{-1}\in M\cap Q=1,
\]
so \(x\) centralizes \(Q\).

Since \(Y\trianglelefteq G\) and \(Q\in \Syl_q(Y)\), the Frattini argument gives
\[
    G=Y\NG{G}{Q}.
\]
Because \(Y=MQ\) and \(Q\leq \NG{G}{Q}\), this becomes
\[
    G=M\NG{G}{Q}.
\]
Moreover,
\[
    M\cap \NG{G}{Q}=\NG{M}{Q}=1.
\]
Therefore \(\NG{G}{Q}\) is a complement to \(M\) in \(G\).
\end{proof}

\begin{theorem}[Hall's theorem: existence criterion]
Let \(G\) be a finite group, and write
\[
    |G|=p_1^{a_1}p_2^{a_2}\cdots p_t^{a_t},
\]
where \(p_1,\ldots,p_t\) are distinct primes. Then \(G\) is solvable if and only if
\(G\) has a Hall \(\pi\)-subgroup for every subset
\[
    \pi\subseteq \pi(G).
\]
\end{theorem}

\begin{proof}
We prove both directions.

\medskip

\noindent
\textbf{First direction: assume \(G\) has Hall \(\pi\)-subgroups for every
\(\pi\subseteq \pi(G)\). We prove that \(G\) is solvable.}

We use induction on \(|G|\). If \(|\pi(G)|=1\), then \(G\) is a \(p\)-group, hence
solvable. If \(|\pi(G)|=2\), then
\[
    |G|=p^a q^b
\]
for two primes \(p,q\), so \(G\) is solvable by Burnside's \(p^a q^b\)-theorem.

Now assume \(|\pi(G)|=t\geq 3\). For \(1\leq i\leq t\), let
\[
    H_i
\]
be a Hall \(p_i'\)-subgroup of \(G\), where
\[
    p_i'=\pi(G)\setminus \{p_i\}.
\]

We first note that each \(H_i\) is solvable. Indeed, \(H_i\) inherits Hall
subgroups as follows. If \(\sigma\subseteq \pi(H_i)\), take a Hall
\(\sigma\cup\{p_i\}\)-subgroup \(L\) of \(G\). Since
\[
    [G:H_i]
\]
is a power of \(p_i\), while
\[
    [G:L]
\]
is a \((\sigma\cup\{p_i\})'\)-number, these two indices are coprime. Hence
\[
    G=H_iL.
\]
Therefore
\[
    |H_i\cap L|
    =
    \frac{|H_i||L|}{|G|},
\]
and this is precisely the \(\sigma\)-part of \(|H_i|\). Thus \(H_i\cap L\) is a
Hall \(\sigma\)-subgroup of \(H_i\). By induction, \(H_i\) is solvable.

In particular, \(H_1\) is solvable. Hence \(H_1\) contains a nontrivial minimal
normal subgroup, so there exists a prime \(p\in \pi(H_1)\) such that
\[
    \Op{p}{H_1}\neq 1.
\]
After relabelling the primes if necessary, we may assume
\[
    p\in \{p_3,\ldots,p_t\}.
\]
Thus \(p\mid |H_2|\). Let \(P\in \Syl_p(H_2)\). Then \(P\) is also a Sylow
\(p\)-subgroup of \(G\). By Sylow's theorem, there exists \(g\in G\) such that
\[
    \Op{p}{H_1}\leq P^g\leq H_2^g.
\]

Now
\[
    |H_1|
    =
    p_2^{a_2}p_3^{a_3}\cdots p_t^{a_t},
\]
and
\[
    |H_2^g|
    =
    p_1^{a_1}p_3^{a_3}\cdots p_t^{a_t}.
\]
Moreover,
\[
    |H_1\cap H_2^g|
    \leq
    p_3^{a_3}\cdots p_t^{a_t}.
\]
Therefore
\[
\begin{aligned}
    |H_1H_2^g|
    &=
    \frac{|H_1||H_2^g|}{|H_1\cap H_2^g|}  \\
    &\geq
    \frac{
        (p_2^{a_2}p_3^{a_3}\cdots p_t^{a_t})
        (p_1^{a_1}p_3^{a_3}\cdots p_t^{a_t})
    }{
        p_3^{a_3}\cdots p_t^{a_t}
    }  \\
    &=
    |G|.
\end{aligned}
\]
Since \(H_1H_2^g\subseteq G\), we must have
\[
    G=H_1H_2^g.
\]

Let
\[
    N=\left\langle \Op{p}{H_1}^{\,x}\mid x\in G\right\rangle
\]
be the normal closure of \(\Op{p}{H_1}\) in \(G\). Then
\[
    1\neq N\trianglelefteq G.
\]
We claim that \(N<G\). Since \(G=H_1H_2^g\), every \(x\in G\) may be written as
\[
    x=ab,
    \qquad
    a\in H_1,\ b\in H_2^g.
\]
Because \(\Op{p}{H_1}\trianglelefteq H_1\),
\[
    \Op{p}{H_1}^{\,x}
    =
    \Op{p}{H_1}^{\,ab}
    =
    \Op{p}{H_1}^{\,b}
    \leq H_2^g.
\]
Hence all conjugates of \(\Op{p}{H_1}\) lie inside \(H_2^g\), and therefore
\[
    N\leq H_2^g<G.
\]
Thus \(N\) is a proper nontrivial normal subgroup of \(G\).

By the previous lemma, normal subgroups and quotient groups inherit the Hall
existence property from \(G\). Hence both \(N\) and \(G/N\) have Hall
\(\pi\)-subgroups for every relevant set of primes. By induction, \(N\) and
\(G/N\) are solvable. Therefore \(G\) is solvable.

\medskip

\noindent
\textbf{Second direction: assume \(G\) is solvable. We prove that all Hall
\(\pi\)-subgroups exist.}

Again we use induction on \(|G|\). Let
\[
    M\trianglelefteq G
\]
be a minimal normal subgroup. Since \(G\) is solvable,
\[
    M\cong (C_p)^d
\]
for some prime \(p\) and some \(d\geq 1\). Put
\[
    \overline{G}=G/M.
\]
By induction, \(\overline{G}\) has Hall subgroups for every subset of
\(\pi(\overline{G})\).

Fix \(\pi\subseteq \pi(G)\).

\medskip

\noindent
\textbf{Case 1: \(p\in \pi\).}

Let \(\overline{H}\) be a Hall
\[
    \pi\cap \pi(\overline{G})
\]
subgroup of \(\overline{G}\). Let \(X\) be the full preimage of
\(\overline{H}\) under the quotient map \(G\to \overline{G}\). Then
\[
    M\leq X\leq G
\]
and
\[
    X/M=\overline{H}.
\]
Since \(M\) is a \(p\)-group and \(p\in \pi\), the subgroup \(X\) contains the full
\(p\)-part of \(G\), and \(\overline{H}\) contains the full \(\pi\)-part of
\(\overline{G}\). Hence \(X\) is a Hall \(\pi\)-subgroup of \(G\).

\medskip

\noindent
\textbf{Case 2: \(p\notin \pi\).}

Let \(\overline{H}\) be a Hall \(\pi\)-subgroup of \(\overline{G}\), and let \(X\)
be its full preimage in \(G\). Then
\[
    X/M=\overline{H},
\]
and \(M\) is a normal \(p\)-subgroup of \(X\), with \(p\nmid |\overline{H}|\).

If \(X<G\), then \(X\) is solvable and has smaller order than \(G\). By induction,
\(X\) has a Hall \(\pi\)-subgroup \(H\). Since \(X/M=\overline{H}\) contains the
full \(\pi\)-part of \(G/M\), this \(H\) is also a Hall \(\pi\)-subgroup of \(G\).

If \(X=G\), then
\[
    G/M=\overline{H}
\]
is a \(\pi\)-group and \(p\nmid |G/M|\). By the complement step, \(M\) has a
complement \(H\) in \(G\). Thus
\[
    G=MH,
    \qquad
    M\cap H=1,
    \qquad
    H\cong G/M.
\]
Since \(G/M\) is a \(\pi\)-group and \(M\) is a \(p\)-group with \(p\notin \pi\),
this complement \(H\) is a Hall \(\pi\)-subgroup of \(G\).

Thus, in every case, \(G\) has a Hall \(\pi\)-subgroup.
\end{proof}

\section{Conjugacy of Hall Subgroups in Solvable Groups}

\begin{proposition}[Conjugacy of Hall subgroups]
Let \(G\) be a finite solvable group, and let \(\pi\subseteq \pi(G)\). Then all
Hall \(\pi\)-subgroups of \(G\) are conjugate in \(G\).
\end{proposition}

\begin{proof}
We use induction on \(|G|\). Let \(H\) and \(K\) be Hall \(\pi\)-subgroups of
\(G\).

Since \(G\) is solvable, there exists a maximal normal subgroup
\[
    M\trianglelefteq G
\]
such that
\[
    |G/M|=p
\]
for some prime \(p\). Equivalently,
\[
    G/M\cong C_p.
\]

\medskip

\noindent
\textbf{Case 1: \(p\notin \pi\).}

Then \(H\leq M\) and \(K\leq M\). Indeed, if \(H\not\leq M\), then
\[
    HM=G,
\]
so
\[
    H/(H\cap M)\cong G/M,
\]
and hence \(p\mid |H|\), contradicting \(p\notin \pi\).

Thus \(H\) and \(K\) are Hall \(\pi\)-subgroups of the solvable group \(M\).
By induction, they are conjugate in \(M\), hence also conjugate in \(G\).

\medskip

\noindent
\textbf{Case 2: \(p\in \pi\).}

Then
\[
    HM=G
    \qquad\text{and}\qquad
    KM=G.
\]
Moreover,
\[
    H\cap M
    \qquad\text{and}\qquad
    K\cap M
\]
are Hall \(\pi\)-subgroups of \(M\). By induction, there exists \(x\in M\) such
that
\[
    (H\cap M)^x=K\cap M.
\]
Replacing \(H\) by \(H^x\), we may assume
\[
    H\cap M=K\cap M.
\]
Set
\[
    L=H\cap M=K\cap M.
\]
Then \(L\trianglelefteq H\) and \(L\trianglelefteq K\), because \(M\trianglelefteq G\).
Hence
\[
    H,K\leq \NG{G}{L}.
\]

If
\[
    \NG{G}{L}<G,
\]
then \(H\) and \(K\) are Hall \(\pi\)-subgroups of the solvable group
\(\NG{G}{L}\). By induction, \(H\) and \(K\) are conjugate in \(\NG{G}{L}\), hence
in \(G\).

It remains to consider the case
\[
    \NG{G}{L}=G.
\]
Then \(L\trianglelefteq G\). Passing to the quotient \(G/L\), we have
\[
    H/L\cong G/M\cong C_p
    \qquad\text{and}\qquad
    K/L\cong G/M\cong C_p.
\]
Thus \(H/L\) and \(K/L\) are Sylow \(p\)-subgroups of \(G/L\). By Sylow's
conjugacy theorem, there exists \(g\in G\) such that
\[
    (H/L)^g=K/L.
\]
Since \(L\leq H,K\), this gives
\[
    H^g=K.
\]
Therefore \(H\) and \(K\) are conjugate in \(G\).
\end{proof}

\begin{remark}
For finite solvable groups, Hall subgroups behave like Sylow subgroups in two
essential ways:
\begin{enumerate}[label=\textup{(\arabic*)}]
    \item Hall \(\pi\)-subgroups exist for every \(\pi\subseteq \pi(G)\);
    \item any two Hall \(\pi\)-subgroups are conjugate.
\end{enumerate}
Both statements can fail for nonsolvable groups.
\end{remark}